% Nová příprava: 2026-10-28; 1. ročník SPŠ, etapa 13–24.
% Rozpis: 1_rocnik/tematicke_plany/Rozpis_vyuky_MAT_1_rocnik.tex, hodina 17.
% Původ: vlastní formulace výkladu a vlastní úlohy; metodická inspirace uvedenými zdroji.
% Kontrola, obsahové porovnání zdrojů a časy: planovani/DOKONCENI_PRIPRAV_13_24.md.
% Učebnice: ucebnice/01_Zakladni_poznatky/Mnoziny/Množiny.tex
% oddíly Rovnost množin, Podmnožiny, Množinové operace a Vennovy diagramy; dle tématu hodiny.
% Realisticky.cz: externi_zdroje/realisticky/01_Zakladni_poznatky/03_Mnoziny/010303_Mnozinove_operace/010303_Lekce.pdf
% Realisticky.cz: externi_zdroje/realisticky/01_Zakladni_poznatky/03_Mnoziny/010304_Vennovy_diagramy/010304_Lekce.pdf
% Realisticky.cz: externi_zdroje/realisticky/01_Zakladni_poznatky/03_Mnoziny/010309_Mnoziny_shrnuti/010309_Lekce.pdf
\documentclass[11pt,a4paper]{article}
% Spolecny styl pro ucitelske pripravy MAT (Work/MAT)
% Verze 1.1, 2026-10-23
\NeedsTeXFormat{LaTeX2e}
\ProvidesPackage{priprava}[2026/10/23 v1.1 Sablona ucitelskych priprav MAT]

% Zakladni typografie a cestina
\RequirePackage[T1]{fontenc}
\RequirePackage[utf8]{inputenc}
\RequirePackage[czech]{babel}
\RequirePackage{lmodern}
\RequirePackage{microtype}
\RequirePackage{amsmath,amssymb}
\RequirePackage{xcolor}
\RequirePackage{enumitem}
\RequirePackage[most]{tcolorbox}
\RequirePackage{tabularx,booktabs,array}
\RequirePackage{fancyhdr}
\RequirePackage{needspace}
\RequirePackage{geometry}
\RequirePackage[hidelinks]{hyperref}

\geometry{a4paper,top=17mm,bottom=18mm,left=23mm,right=23mm}
\definecolor{notegray}{gray}{0.39}
\setlength{\parindent}{0pt}
\setlength{\parskip}{2.5pt plus 1pt minus 0.5pt}
\setlist[itemize]{topsep=2pt,itemsep=1pt,parsep=0pt,leftmargin=1.45em}
\setlist[enumerate]{topsep=2pt,itemsep=2pt,parsep=0pt,leftmargin=1.65em}
\setcounter{secnumdepth}{0}
\renewcommand{\arraystretch}{1.12}

% Zahlavi: vlevo rocnik | globalni tema, vpravo cislo: nazev hodiny
% Zapati zustava zcela prazdne.
\pagestyle{fancy}
\fancyhf{}
\renewcommand{\headrulewidth}{0.35pt}
\renewcommand{\footrulewidth}{0pt}
\setlength{\headheight}{14pt}
\newcommand{\nastavPripravu}[4]{%
  \fancyhead[L]{\small #1\enspace\textbar\enspace #2}%
  \fancyhead[R]{\small #3: #4}%
}

% Minimalisticke cernobile boxy. Kazda definice zvlast.
\tcbset{commonlesson/.style={enhanced,breakable=false,colback=white,colframe=black,
  colbacktitle=white,coltitle=black,boxrule=0.45pt,
  arc=0.8mm,outer arc=0.8mm,left=3mm,right=3mm,top=1.6mm,bottom=1.6mm,
  toptitle=1mm,bottomtitle=0.7mm,fonttitle=\bfseries,
  attach boxed title to top left={xshift=2.5mm,yshift=-1.15mm},
  boxed title style={colback=white,colframe=white,boxrule=0pt,left=1mm,right=1mm}}}
\newtcolorbox{zadaniBox}[1][]{commonlesson,title={Zadání},#1}
\newtcolorbox{prikladBox}[1][]{commonlesson,title={Příklad},#1}
\newtcolorbox{definiceBox}[1][]{commonlesson,title={Definice},#1}

% Spolecny box pro zadani a jeho reseni, bez mezery mezi castmi.
% Pouziti:
%   \begin{prikladReseniBox}
%     Zadani ulohy.
%     \tcblower
%     \begin{align*} ... \end{align*}
%   \end{prikladReseniBox}
% Spodni cast ma jemny sedy podklad a delici linku; lze v ni sazet
% vice radku matematiky. Obe casti maji jen jeden spolecny ramecek.
\tcbset{solutionlesson/.style={
  commonlesson,
  enhanced jigsaw,
  breakable,
  colbacklower=black!3,
  segmentation style={solid,draw=black!25,line width=0.35pt},
  fontlower=\small\color{notegray},
  before lower={\noindent\textbf{Řešení:}\enspace},
  leftlower=3mm,rightlower=3mm
}}
\newtcolorbox{prikladReseniBox}[1][]{solutionlesson,title={Příklad},#1}
\newtcolorbox{zadaniReseniBox}[1][]{solutionlesson,title={Zadání},#1}

% Pomucky pro autora pripravy: sedive poznamky; \reseni pro starsi pripravy
\newcommand{\poznamka}[1]{\par\begingroup\small\color{notegray}\itshape
  \textnormal{\textbf{Poznámka pro učitele.}}\enspace #1\par\endgroup}
\newcommand{\reseni}[1]{\par\begingroup\small\color{notegray}
  \textbf{Očekávané řešení:}\enspace #1\par\endgroup}
\newcommand{\blok}[1]{\needspace{5\baselineskip}\subsection*{#1}}
\newcommand{\vyrok}[1]{\uv{#1}}

\usepackage{tikz}
\usetikzlibrary{patterns,calc,angles,quotes}
\nastavPripravu{1. ročník}{Množiny}{17}{Vlastnosti operací}
\begin{document}
\blok{Pravidlo musí platit pro libovolný prvek}
U čísel jsme zkoumali komutativitu a asociativitu. Také u množin se ptáme,
kdy lze měnit pořadí či závorky. Jeden vypočtený příklad může pravidlo
ilustrovat nebo vyvrátit, ale sám obecnou rovnost nedokazuje.
V celé hodině jsou $A,B,C\subseteq U$ a doplňky tvoříme ve stejném $U$.

\begin{center}
\begin{tabularx}{\linewidth}{@{}l X@{}}
\toprule
Vlastnost & Matematický zápis\\ \midrule
Komutativita & $A\cup B=B\cup A$, \quad $A\cap B=B\cap A$\\
Asociativita & $(A\cup B)\cup C=A\cup(B\cup C)$\\
             & $(A\cap B)\cap C=A\cap(B\cap C)$\\
Distributivita & $A\cap(B\cup C)=(A\cap B)\cup(A\cap C)$\\
              & $A\cup(B\cap C)=(A\cup B)\cap(A\cup C)$\\
Prázdná a základní množina & $A\cup\varnothing=A$, $A\cap U=A$\\
                          & $A\cap\varnothing=\varnothing$, $A\cup U=U$\\
Doplněk & $A\cup A'=U$, $A\cap A'=\varnothing$, $(A')'=A$\\
De Morganovy zákony & $(A\cup B)'=A'\cap B'$,\quad $(A\cap B)'=A'\cup B'$\\
\bottomrule
\end{tabularx}
\end{center}
Komutativita zaměňuje pořadí množin, asociativita jen seskupení stejné operace.
Distributivita propojuje dvě různé operace. Rozdíl komutativní není:
$\{1,2\}\setminus\{2\}=\{1\}$, kdežto $\{2\}\setminus\{1,2\}=\varnothing$.

\begin{prikladReseniBox}[title={P1: konkrétní kontrola distributivity},breakable=false]
Pro $A=\{1,2,4\}$, $B=\{2,3\}$, $C=\{3,4,5\}$ vypočti obě strany
$A\cap(B\cup C)=(A\cap B)\cup(A\cap C)$.
\tcblower
$B\cup C=\{2,3,4,5\}$, takže vlevo $\{2,4\}$.
Vpravo $A\cap B=\{2\}$ a $A\cap C=\{4\}$, tedy také $\{2,4\}$.
Ověřili jsme tuto trojici množin. Obecný důvod potřebujeme zvlášť.
\end{prikladReseniBox}

\begin{prikladReseniBox}[title={P2: proč distributivita platí obecně}]
Vysvětli první distributivní zákon pomocí libovolného prvku $x$ a diagramu.
\tcblower
Prvek $x$ patří vlevo právě tehdy, když je v $A$ a současně alespoň v jedné
z množin $B,C$. To je totéž jako: je v $A$ i $B$, nebo je v $A$ i $C$.
Jde přesně o podmínku příslušnosti k pravé straně.
\[
x\in A\land(x\in B\lor x\in C)
\quad\Longleftrightarrow\quad
(x\in A\land x\in B)\lor(x\in A\land x\in C).
\]
\begin{center}
\begin{tikzpicture}[scale=.66]
\begin{scope}
\clip (-.75,.35) circle(1.35);
\fill[black!20] (.75,.35) circle(1.35);
\fill[black!20] (0,-.7) circle(1.35);
\end{scope}
\draw (-.75,.35) circle(1.35); \draw (.75,.35) circle(1.35); \draw (0,-.7) circle(1.35);
\node at (-1.65,1.65) {$A$}; \node at (1.65,1.65) {$B$}; \node at (0,-2.25) {$C$};
\node at (0,-2.8) {$A\cap(B\cup C)$};
\begin{scope}[xshift=6.2cm]
\begin{scope}
\clip (-.75,.35) circle(1.35); \clip (.75,.35) circle(1.35);
\fill[black!20] (-2.2,-2.1) rectangle (2.2,1.8);
\end{scope}
\begin{scope}
\clip (-.75,.35) circle(1.35); \clip (0,-.7) circle(1.35);
\fill[black!20] (-2.2,-2.1) rectangle (2.2,1.8);
\end{scope}
\draw (-.75,.35) circle(1.35); \draw (.75,.35) circle(1.35); \draw (0,-.7) circle(1.35);
\node at (-1.65,1.65) {$A$}; \node at (1.65,1.65) {$B$}; \node at (0,-2.25) {$C$};
\node at (0,-2.8) {$(A\cap B)\cup(A\cap C)$};
\end{scope}
\end{tikzpicture}
\end{center}
Obě konstrukce vyznačují stejné oblasti včetně společné části všech tří kruhů.
Slovní argument platí pro libovolné $x$, proto dokazuje rovnost množin.
\end{prikladReseniBox}

\begin{prikladReseniBox}[title={P3: negace výběru},breakable=false]
Proč platí $(A\cup B)'=A'\cap B'$? Přelož výsledek do běžné věty.
\tcblower
Prvek $U$ neleží ve sjednocení právě tehdy, když není v $A$ a zároveň není v $B$.
Proto je v obou doplňcích. Například díl, který není ani červený, ani kovový,
patří do doplňku množiny „červené nebo kovové díly“.
Negace „ne alespoň jedna podmínka“ znamená „žádná z obou“, nikoli „právě jedna“.
\end{prikladReseniBox}
\poznamka{Tabulka je přehled, nikoli zadání k opisování. Jádro: významy vlastností,
P1–P3 a S1 a), b), přibližně 38 min práce, organizace a závěr 7 min.
Druhý distributivní zákon stačí přečíst; jeho odvození je rezerva.
S1 c) vybírej podle času nebo dokonči doma.}

\blok{Samostatné procvičení}
\begin{zadaniReseniBox}[title={S1: použít pravidlo a najít protipříklad},breakable=false]
\begin{enumerate}[label=\alph*)]
\item Pro $U=\{1,\ldots,6\}$, $A=\{1,2,4\}$, $B=\{2,3,4\}$
urči $(A\cup B)'$ i $A'\cap B'$ a porovnej.
\item Zdůvodni, proč $A\cup(A\cap B)=A$ pro libovolné $A,B$.
\item Platí vždy $A\setminus(B\cup C)=(A\setminus B)\cup(A\setminus C)$?
Otestuj $A=\{1\},B=\{1\},C=\varnothing$.
\end{enumerate}
\tcblower
a) $A\cup B=\{1,2,3,4\}$, doplněk $\{5,6\}$;
$A'=\{3,5,6\}$ a $B'=\{1,5,6\}$, průnik $\{5,6\}$.
b) $A\cap B\subseteq A$, proto připojením průniku nepřidáme žádný nový prvek;
všechny prvky $A$ zůstávají. Jde o zákon absorpce.
c) Vlevo $\varnothing$, vpravo $\varnothing\cup\{1\}=\{1\}$.
Obecně správně $A\setminus(B\cup C)=(A\setminus B)\cap(A\setminus C)$:
prvek $A$ musí být mimo obě vylučované množiny současně.
\end{zadaniReseniBox}

\blok{Volitelná rezerva}
\begin{zadaniReseniBox}[title={R1: druhé odvození, 10–15 minut},breakable=false]
\begin{enumerate}[label=\alph*)]
\item Vysvětli na libovolném prvku $(A\cap B)'=A'\cup B'$.
\item Dokaž $A\cup(B\cap C)=(A\cup B)\cap(A\cup C)$ rozborem,
zda prvek do $A$ patří, nebo nepatří.
\end{enumerate}
\tcblower
a) Aby prvek nebyl současně v obou množinách, musí alespoň jedné z nich chybět.
Tedy je mimo $A$ nebo mimo $B$, případně mimo obě.
b) Je-li $x\in A$, patří do obou stran bez dalších podmínek.
Není-li $x\in A$, levá strana vyžaduje $x\in B\cap C$.
Pravá strana pak také vyžaduje současně $x\in B$ a $x\in C$.
Případy pokrývají všechny prvky.
\end{zadaniReseniBox}

\blok{Domácí procvičení}
\begin{zadaniReseniBox}[title={D1: 12–18 minut}]
Pro a)–c) je $U=\{1,\ldots,10\}$, $A$ jsou sudá čísla z $U$,
$B=\{1,2,3,4\}$ a $C=\{3,6,9\}$.
\begin{enumerate}[label=\alph*)]
\item Spočti obě strany prvního distributivního zákona.
\item Urči $(A\cup B)'$ a přepočítej výsledek pomocí doplňků.
\item Urči $(A\cap B)'$ a přepočítej výsledek pomocí doplňků.
\item Najdi protipříklad k $A\setminus B=B\setminus A$.
\item Zjednoduš $M\cap(M\cup N)$ a výsledek zdůvodni pro libovolný prvek.
\end{enumerate}
\tcblower
a) Obě strany $\{2,4,6\}$: $B\cup C=\{1,2,3,4,6,9\}$,
$A\cap B=\{2,4\}$, $A\cap C=\{6\}$.
b) $\{5,7,9\}$; $A'=\{1,3,5,7,9\}$, $B'=\{5,6,7,8,9,10\}$,
jejich průnik je uvedený výsledek.
c) $\{1,3,5,6,7,8,9,10\}$; jde o sjednocení týchž dvou doplňků.
d) Např. $A=\{1\}, B=\{2\}$, rozdíly $\{1\}$ a $\{2\}$.
e) $M$: prvek $M$ automaticky splňuje i druhou podmínku;
prvek mimo $M$ nemůže být v průniku.
\end{zadaniReseniBox}
\end{document}
